% Laget automatisk av elevark/lag_elevark.py fra opplegg/elektrisk-felt-vannbad.html. Endringer her blir overskrevet.
\documentclass[11pt]{article}
\usepackage{elevark}
\begin{document}
\elevarkhode{Fysikk 2 · Elektriske felt og krefter}{Elektrisk felt i et vannbad}{https://fysikklab.no/opplegg/elektrisk-felt-vannbad.html}{Opplegget på nett}

Et elektrisk felt kan ikke ses, men potensialet kan måles. I et grunt vannbad med spenning mellom to elektroder kan dere måle potensialet punkt for punkt med et voltmeter, tegne ekvipotensiallinjer og finne feltlinjene som står vinkelrett på dem.

\section*{Hensikt}

\begin{itemize}
\item Kartlegge ekvipotensiallinjer og tegne feltlinjer for to ulike elektrodeoppsett.
\item Vise at feltet mellom to parallelle plater er tilnærmet homogent, med $E = U/d$.
\end{itemize}

\section*{Utstyr}

\begin{itemize}
\item Et grunt plastbrett med omtrent 1~cm vanlig springvann, og millimeterpapir under brettet (i en plastlomme).
\item To rette metallstriper (for eksempel kobberbånd) som elektroder i del A, og to metallstaver eller skruer i del B.
\item Spenningskilde på omtrent 6~V (vekselspenning gir minst problemer med gassbobler ved elektrodene), ledninger og et multimeter med en spiss målepinne.
\end{itemize}

\section*{Teorien}

Multimeteret måler potensialforskjellen mellom minuspolen og punktet der målepinnen står. Linjer med samme potensial kalles ekvipotensiallinjer. Feltlinjene står vinkelrett på dem og peker mot lavere potensial. Feltstyrken er hvor fort potensialet endrer seg: $E = ΔV/Δx$. Mellom to parallelle plater med spenningen $U$ og avstanden $d$ er feltet tilnærmet homogent, og $E = U/d$.

\section*{Del A: To parallelle plater}

\begin{enumerate}
\item Legg de to stripene parallelt i vannet med avstanden $d = 20\,\mathrm{cm}$. Koble dem til spenningskilden, og mål spenningen $U$ mellom dem.
\item Koble multimeterets COM til minuspolen. Mål potensialet hver andre centimeter langs en linje rett fra minuspolen til plusspolen, og skriv det inn i tabellen.
\item Mål også et par punkter ved siden av linjen. Har punkter med samme avstand fra minuspolen samme potensial?
\end{enumerate}

\section*{Del B: To punktelektroder}

\begin{enumerate}
\item Bytt stripene med to staver som står loddrett i vannet, omtrent 15~cm fra hverandre.
\item Finn punkter med potensial 1~V, 2~V, 3~V, 4~V og 5~V. Flytt målepinnen til multimeteret viser den samme verdien, og merk punktene på et eget ark med samme rutenett.
\item Trekk ekvipotensiallinjene gjennom punktene. Tegn feltlinjer som står vinkelrett på dem, fra plusspolen til minuspolen.
\end{enumerate}

\section*{Oppgaver}

\begin{oppgaveliste}
\oppgave \fase{gjett}{Gjett}Hvordan tror dere potensialet endrer seg fra minuspolen til plusspolen i del A? Jevnt, eller raskest nær en av platene?
\svar
\oppgave \fase{test}{Mål}Gjør del A, eller trykk på «Fyll inn eksempel». Ligger punktene på en rett linje? Sammenlign $E$ fra stigningstallet med $U/d$.
\svar
\oppgave \fase{forklar}{Forklar}Hvorfor står feltlinjene vinkelrett på ekvipotensiallinjene? Hvor mye arbeid gjør feltet på en ladning som flyttes langs en ekvipotensiallinje?
\svar
\oppgave \fase{test}{Tegn}Gjør del B. Hvor ligger ekvipotensiallinjene tettest, og hva betyr det for feltstyrken der? Sammenlign tegningen med simuleringen for to motsatte ladninger.
\svar
\oppgave \fase{forklar}{Vurder}Linjen i del A går ikke gjennom 0~V og $U$ ved platene. Hva kan være grunnen?
\svar
\oppgave \fase{utfordring}{Utfordring}Et elektron befinner seg i et homogent felt på 27~V/m. Hvor stor kraft virker på det, og hvor stor akselerasjon får det? Hvorfor er dette forsøket likevel ikke farlig?
\svar
\end{oppgaveliste}

\end{document}
